\documentclass[11pt]{article}
\usepackage{amsmath,amssymb,amsthm,hyperref}
\newtheorem{proposition}{Proposition}
\begin{document}
\title{Adjacent-grid Cauchy profiles and their first support obstruction}
\date{November 24, 2026}
\maketitle

\paragraph{Status.}
This is a discarded global construction, with a still meaningful auxiliary
local decoding question. It proves no improvement in the worst-case head
bound. The original model and admissibility conditions are unchanged.
The three elementary finite support statements below are kernel checked
in \texttt{LargeDomainGridSupport.lean}. The rational construction and its
interpretation in the Boolean cube are written arguments.

\section{The auxiliary decoder}
There are $a$ row blocks, each with $m$ one-hot coordinates. Source cells
are $u\in\{0,\ldots,m-1\}^a$ and vertices are
$v\in\{0,\ldots,m\}^a$. The nonnegative affine coarse coefficient at
row coordinate $(i,t)$ is
\[
 C_{v,i,t}=(t+1/2-v_i)^2-1/4\geq0.
\]
It is zero exactly when $v_i=t$ or $v_i=t+1$. Thus on the valid one-hot
row input $u$, precisely its $2^a$ adjacent vertices have zero coarse
penalty. There are $G=m^a$ source cells and $H=(m+1)^a$ vertices. For
$m\geq a$, $H/G$ is bounded by an absolute constant. Every row/query
input can therefore interact with exponentially many heads in $a$;
this passes the sparse-interaction degree screen that defeats polynomial
active overlap.

A generic nonnegative affine fine penalty $p_v$ can give all $2^a$
adjacent cells distinct levels. Consider the local score
\[
 F_j(u,i)=\sum_{v\text{ adjacent to }u}
 \frac{c_{vj}q_{vi}}{(b_{vj}+p_v(u))^2},
 \qquad 1<b_{vj}<1+\eta.
\]
The vectors $q_v\in\mathbb R^a$ may be arbitrary and shared across the
independently labeled columns $j$. The amplitudes $c_{vj}$ and poles
$b_{vj}$ vary independently with the column. All of them may be chosen
jointly after seeing the entire label family. This is not forced to be a
gradient field: its query vectors need not equal corner displacements.

Two heads per vertex give the displayed squared reciprocal as a query
slope tangent. Both heads have the same coarse gate, the same column
poles, opposite column numerators, and query slopes differing by a small
multiple of $q_v$. A common slope guard permits either sign of each
$q_{vi}$. Their score is exactly zero on the query axis and the column
axis. Complementing all source coordinates makes the row slopes
negative, consistently with the reverse local interfaces; taking the
column embedding intercept above all compact poles also gives negative
column slopes and a positive full-cube floor. The query tangent is taken
first and the coarse penalty tends to infinity afterwards. These facts
only provide a local score, not the required support on earlier layers.

\section{The unique all-double input}
For an interior vertex, $1\leq v_i\leq m-1$ for every $i$, select both
row coordinates $v_i-1$ and $v_i$ in every block. Call this forward input
$X_v$. It has weight $2a$. Its complement has reverse row weight
$am-2a$, strictly below the target reverse row weight $am-a$ when $a>0$.
A vertex $w$ has zero coarse penalty at $X_v$ if and only if, for every
$i$, it is adjacent to both intervals $v_i-1$ and $v_i$. This forces
$w_i=v_i$. Therefore only the vertex $v$ survives.

\begin{proposition}
At $X_v$, the limiting score on query $i$ and column $j$ is exactly
\[
 \frac{c_{vj}q_{vi}}{(b_{vj}+p_v(X_v))^2}.
\]
If the unmodified gated family must vanish on all earlier reverse layers,
then $c_{vj}q_{vi}=0$ for every interior $v$, every query $i$, and every
column $j$.
\end{proposition}
\begin{proof}
The denominator is strictly positive on the full cube. The unique-gate
observation removes every other summand, so a zero score forces the
numerator product to be zero. In Lean, the declarations are
\texttt{grid\_double\_compatible\_iff},
\texttt{grid\_double\_score\_eq}, and
\texttt{grid\_double\_zero\_forces\_residue}, in namespace
\texttt{HeadComplexity.LargeDomain}.
\end{proof}

This rules out using the unmodified auxiliary decoder in the current
reverse graded-shell compiler. It does not rule out extra cancellation
banks, different source domains, different lower-layer behavior, or a
nontriangular global construction. In particular, it is not a lower
bound for unrestricted original-model heads.

The tempting count $((m+1)^a-m^a)$ of missing-block exceptions uses the
opposite orientation and is not valid here. After the required source
complement, earlier reverse layers include multihot forward blocks.
A formal product with roots at all row weights $a+1,\ldots,2a$ would
remove this obstruction, but its degree is $a$ and there is no established
cheap shared implementation that retains the generic fine Cauchy levels.

\section{Finite exploratory calculation}
A PyTorch search used $a=4,m=8,s=16$, so $G=4096,H=6561$, and
$262144$ independent random signs. With fixed generic fine slopes and
compact pole width $0.05$, the trainable parameters were the shared
$q_{vi}$, the column amplitudes and offsets, and cheap query/column
biases. There were $236260$ real parameters, fewer than target entries.
After $22000$ Adam steps, $230$ sign constraints remained unsatisfied.
This is only an optimization failure. It proves neither local
nonrepresentability nor a positive exact decoder theorem. The support
counterexample above is independent of this computation, and is the
reason this route was stopped as a global candidate.

\end{document}
